\section{Discussion}

\subsection{Analysis of Interviews with Performing Musicians}

Interviews with the performing musicians were conducted individually: with Mr. Yoshijima on December 15, 2017, with Mr. Kawamoto on December 18 of the same year, and with Mr. Kaneko on December 19 of the same year. The interviews were conducted as semi-structured interviews with predetermined topics aimed at investigating the musicians' subjective experiences regarding the following items, while prioritizing the natural flow of conversation and allowing free expression. The full interview content is provided in \ref{appendix}, so excerpts will be discussed here. The items are as follows:

\begin{enumerate}
  \item Impressions after listening to the video and audio recordings of this work
  \item Regarding the work's booklet, was the content understandable?
	\item Impressions of the booklet and whether there was any useful knowledge gained?
  \item Have you experienced any similar music in the past?
	\item Are there any differences compared to what you usually do?
	\item Where do you think the difficulty of performance lies?
  \item What do you think about the potential development of this work and method?
\end{enumerate}

\subsubsection{Regarding "Impressions after listening to the video and audio recordings of this work"}

Both Yoshijima and Kawamoto expressed concerns about whether actual performance was possible. This can be understood from Yoshijima's comment that "there is no precedent for such a score" and Kawamoto's remark that "there is content written in the score, but I wondered if this would be okay." However, both felt a gap between the score and the music, noting "it became music unexpectedly" and "it became more proper music than I thought." Kaneko also gave positive evaluation regarding the feasibility of performance. Concerning the score and actual performance, Yoshijima pointed out: "I don't know what kind of sound it will become when the three of us come together, or what kind of groove will emerge." Kawamoto also raised the point "I understood the concept, but I wondered if it wasn't impossible to perform," suggesting that he felt the difficulty from the work itself.

\subsubsection{Regarding "Was the content of the work's booklet understandable?"}

While Yoshijima and Kawamoto stated they "understood," Kaneko said "I haven't read it thoroughly enough, so it hasn't sunk in yet." As Kaneko mentioned in the interview that he was "approaching it as business," it appears he was viewing it from the standpoint of an ordinary "job" and focused only on what was necessary to complete the commission.

\subsubsection{Regarding "Impressions of the booklet and whether there was any useful knowledge gained?"}

Yoshijima, the drummer, expressed this attempt as "like having one rail, and although there should be another flowing in a different direction, somehow they're on the same rail," and stated that "I can already apply this within myself and bring it out in other situations," suggesting that understanding has progressed to the level of physical embodiment. On the other hand, Kawamoto noted "from a musician's musical information perspective, there are parts that are slightly over-explained. It could be simpler," while also stating "but the way the drum score is divided into two staves and such ingenuity is very good (omitted). It was very good for understanding the layers of music. The fact that it shows how this and this add up to become this was very good. As a score that appears in the field, it's unlike anything I've seen before, but I think it's sufficient as an explanation of the project," acknowledging a certain effectiveness of the booklet for performing this work.

\subsubsection{Regarding "Have you experienced any similar music in the past?"}

We were able to hear about many examples of polyrhythm/cross-rhythm in general and experiences of attempting them personally, but similarities to the characteristics of this work—such as the way rhythms are taken and the notation—remained at the level of "possible that similar things existed in the past." This is due to the impossibility of confirming all music, and comments were also heard regarding the possibility of analysis arising from the notation method.

\subsubsection{Regarding "Are there any differences compared to what you usually do?"}

Yoshijima's answer that "what I usually do is more difficult" is interesting. The reason being "when each person is only doing their own thing, even if I try to connect, it's impossible to hold them together," and compared to such situations, he wondered "is it easier when the performers are trying to connect?" When asked this, he responded: "That's really just a matter of whether it's musical. In that sense, for this piece, I felt more comfortable. No matter how high an individual's ability is, it's more painful in a field where people gather who perform selfishly." I speculate that this is also influenced by the fact that the drummer's role had a large proportion in connecting the parts. Additionally, Kawamoto expressed understanding of this research's positioning as "research on vocabulary," stating "I think it's research on grammar focused on one point, and it's something that really should be done, but normally it's difficult to focus that narrowly and spend time on just that one thing. That's why even though I understood the concept, I didn't know if I could perform it. That's why it's something very valuable," commenting on the research's value in comparison to his usual activities.

\subsubsection{Regarding "Where do you think the difficulty of performance lies?"}

Kaneko cited "because the rhythms are compound," pointing to the basic concept of this work. Regarding the difficulty, while stating "I've heard that when humans do two things, they can't do them at 50\% each but rather at about 30\%, making it inefficient. I think it's that kind of difficulty," he actually considered that "when humans process something like this time, I don't think they're doing it multitask. Isn't it similar to the drummer not processing multitask?" Yoshijima mentioned regarding solos: "when someone takes a solo, someone should be the foundation, but this time the drum had a large proportion of that, though the difficulty changes depending on who produces that pulse." As an example, he stated: "for instance, if the drums start taking a solo, to whom do we entrust the rhythm? I would probably destroy it with considerable force, so by doing that, I think the difficulty level increases further," suggesting that he considers the difficulty of performance to vary depending on relationships. I believe this is due to qualitative differences in performance arising from the width created by the piece being meta-composition and by performers' interpretations. Furthermore, Kawamoto pointed out regarding physicality, from the interpretation of "research on vocabulary" mentioned earlier, that "even if you learn a new word, to be able to use that word freely, you need to internalize it physically once," indicating that this takes time. However, regarding the amount of time, he showed that the proportion of usual experience is large: "I've never thought about 3 over 4 with the resolution used this time. In other words, I could play it because I normally listened, learned, and played repeatedly, so if I had normally only listened to and learned 5 over 4 songs, I probably could do it." This also connects with both Kaneko and Kawamoto expressing during rehearsal that the difficulty of "5 over 4" was due to not having opportunities to play it in regular performance.

\subsubsection{Regarding "What do you think about the potential development of this work and method?"}

Yoshijima cited that understanding of what he usually does in his activities progressed from the reproducibility that the score possesses. This means it is directly influencing Yoshijima's ongoing activities, suggesting the possibility that Yoshijima himself will develop/inherit it. Kawamoto cited examples of counterpoint and Bartók's quartets, suggesting that he recognizes its applicability more as a compositional method than improvisation.

\subsection{Analysis from Interviews with Musicians of Other Genres}

Interviews were also conducted with musicians of other genres who did not participate in the performance. The interview was conducted simultaneously with both Mr. Shinya Ando, a santoor player who is a performer of Indian classical music, and Mr. Jun Haraguchi, a tabla player. The interview items are basically the same content as the previous interview, except for some parts. The interview content is provided in \ref{appendix}, so excerpts will be discussed here.

% \begin{enumerate}
%   \item Impressions after listening to the video and audio recordings of this work
%   \item Regarding the work's booklet, was the content understandable?
% 	\item Impressions of the booklet and whether there was any useful knowledge gained?
%   \item Have you experienced any similar music in the past?
% 	\item Are there any differences compared to what you usually do?
% 	\item Where do you think the difficulty of performance lies?
%   \item What do you think about the potential development of this work and method?
% \end{enumerate}

\subsubsection{Regarding "Impressions after listening to the video and audio recordings of this work"}

Positive comments were observed regarding the characteristics of the work, such as rhythmic changes at the rhythm-switching parts of Etude \#2 and being able to take either beat. Additionally, Ando mentioned the difficulty as a solo performer, relating it to his own experience, and pointed out that to make it possible, "first everyone needs to be familiar with irregular time signatures, with an era when songs in 5 are normally abundant, and when that era comes, it will gradually become possible to enter suddenly with different rhythms." This indicates that in Ando's subjective view, the standard rhythm in the current music situation differs from this piece and Ando's usual experience.

\subsubsection{Regarding "Was the content of the work's booklet understandable?"}

While stating that understanding was possible by examining both the audio source and the score, Ando noted: "Since I've talked with the composer, I have that preconception, so I don't know if I could understand it if I saw this for the first time without knowing anything," reserving judgment about understanding the booklet at first sight. The author also understands that understanding the booklet requires a certain level of musical literacy, and considers this a valid point.

\subsubsection{Regarding "Impressions of the booklet and whether there was any useful knowledge gained?"}

Regarding the possibility of a performer playing two rhythms simultaneously through this research's attempt and what they are currently doing: "Since instruments like santoor and tabla can use both hands, there are various examples of two performers each playing different rhythms, but there have been no past approaches where one performer plays different beats with both hands, not in the way the drummer here is doing. Ando said in the summer, after meeting Yamato, that things he had been thinking about vaguely took form, and he's actually been doing that recently." This is similar to Yoshijima's understanding and application in the field mentioned earlier. I also speculate that being a percussionist may be related to applicability.

\subsubsection{Regarding "Have you experienced any similar music in the past?"}

Haraguchi cited an example performed by Zakir Hussain and Mickey Hart in 1976, "though it doesn't exactly correspond to this research's example." This is "an example of creating one melody within 102.5 beats according to the song's explanation," taking rhythm in a very large cycle. However, Ando stated: "What Zakir Hussain is doing is arranging and changing pieces, but the way of thinking is only Indian music. Even in Shakti and other things, he changes various things but only does Indian classical music," noting that the basic structure of the music is considered within the structure of Indian classical music.

\subsubsection{Regarding "Where do you think the difficulty of performance lies?"}

Ando stated that such attempts themselves are not difficult if you create something like the rhythm hit point tables shown in this paper, but cited "the time it takes to master." Additionally, noting that the difficulty differs between ensemble and individual, he explained: "When you try to do it alone. It becomes a bit easier to do it as a group, because once you learn how to fit it in, once it enters your own system, if you have in yourself how it goes in 5, how it goes in 7, how 3 goes, how 6 and 12 go, you can just fit them in there one by one, but for one person to start doing different things like this drummer takes time." This is interesting when considered in parallel to the relationship between Yoshijima and Kawamoto/Kaneko in this rehearsal. Yoshijima felt difficulty in showing the hit points of two rhythms alone to connect the bass and piano, and once he could do that, he became able to hear the entire piece. Kawamoto and Kaneko performed by learning what Ando calls "how to fit in," or in other words, the relationships. As Kawamoto mentioned in the interview, there was difficulty because it wasn't in their vocabulary, but as Ando points out, having such vocabulary within oneself makes it easier, which matches what Kawamoto himself mentioned as familiarity with 5 over 4.

Additionally, comments about how to listen to each other in ensemble are interesting. Regarding when Ando and Haraguchi perform polyrhythmic things, "When doing polyrhythmic tricks more than usual, the hit points must align so we're listening considerably. The side doing the original beat must not get pulled in, so while listening to the other, they do it like a metronome, but the side approaching is listening more, so to speak. They fit in phrases while listening," revealing differences in states of consciousness regarding hit points in polyrhythmic performance.

\subsubsection{Regarding "What do you think about the potential development of this work and method?"}

While Haraguchi relatively concretely mentioned his own performance, stating "there are elements of this that can be incorporated into my performance," both noted "if you analyze it carefully, you can't understand unless you really face it," but "if you listen to it as music, music that's really beautiful or something like that should spread more," which, regardless of whether this work is beautiful, matches the purpose of "at first glance you can listen normally, but structurally there's something complex behind it. Considering future music development in the context of popular music."

\subsection{Discussion on Composition and Production Environment}

This research, borrowing Kawamoto's words, focuses on rhythm within compositional methods as "research on vocabulary" in music. When viewing composition broadly through the works from there, as shown in \ref{Western-Modern-AmericanNewMusic}, if the compositional freedom before classical notation was established in the 14th and 15th centuries had remained in the present day, such music might have been possible. Additionally, the possibility that similar examples might exist in existing music due to what becomes invisible/visible through notation methods can also be said from the interviews. When considering composition from there, if current compositional methods, whether Western music or popular music, are established by the structural way of thinking about notation that is currently common, then just as Western music could not fully grasp African rhythms, it can be said that possibilities still remain for rhythm as a compositional method.

Furthermore, when considering the act of composition from the act of notation, if scores before lead sheets are analogized to the relationship between "computer programs and compilers," lead sheets can also be understood like the relationship between "structured data and AI." This can also be considered as a difference in physicality, apart from the body in Western music that "reproduces or expresses according to the note values and temperament of the score," as "a body that reproduces structure" in popular music.

Of course, it is self-evident that various musicians are currently taking all kinds of approaches not only to rhythm, so neither is this approach uniquely novel nor are existing approaches becoming obsolete. However, I consider that through the act of composition in popular music in this paper, while considering possibilities for future music to come, it becomes a critical attempt regarding the current state.

If this research is "research on a certain vocabulary," then beyond this research it is possible to develop and inherit toward other vocabularies, or vocabularies that can be considered in relation to this vocabulary. This is evidenced by various styles of popular music that continue to be inherited and developed even today through not only music theories such as from the Schillinger House to the Berklee Method and the Lydian Chromatic Approach in popular music history, but also through bebop, modal, free jazz, noise, and other genres.

In the works created as "music with multilayered temporal structures" this time, we could observe inadequate production environments during the production process. For example, as far as the author investigated, only with Lilypond was it possible to "structurally" describe scores with variable time signature notation. This is also one manifestation that notation presupposes the musical structure currently taken for granted. Additionally, regarding music production software, we could not find software that allows multi-BPM or track settings with multilayered time signatures, and in creating demo recordings, we could only produce in a method that "forced consistency" halfway using Logic Pro X's MIDI stretch function. Of course, it might be possible by making full use of MIDI resolution and placing MIDI data through calculations. However, that can hardly be called a "compositional method from the composer's subjective perspective" as shown with variable time signature notation. I consider that this research has also become a critical attempt regarding existing production environments in that sense.
